\section{Marked spherical cuts}\label{sec:cut}

An embedded round-product comparison collar supplies two pieces of data
for a cut. Its middle sphere marks the boundary identification, and its
retained buffers permit a Yamabe estimate for the two capped classes.
We measure the energy of each entire side in the original metric,
including the part on which the metric is subsequently changed.

\subsection{Marked topology and the energy inequality}

\begin{proposition}[Finite cut criterion]\label{prop:cut-energy}
Let \(g\) be a smooth metric on a closed oriented homotopy four-sphere \(M\). Cut \(M\) along a finite collection of pairwise disjoint smoothly
embedded three-spheres and attach marked standard balls. Give the
capped pieces smooth metrics and the inherited orientations. Write
\[
 m_j=E_+(g;P_j),\qquad E_j=E_+(\widehat g_j),\qquad
 E=E_+(g)=\sum_j m_j .
\]
Suppose the capped metrics have positive Yamabe constants and
\(E_j\leq m_j+\gamma_j\) with \(\gamma_j\geq0\). If
\(\sum_j E_j\leq2T\), in particular if \(E+\sum_j\gamma_j\leq2T\),
at most one capped piece has \(E_j>T\). The original manifold is
oriented-diffeomorphic to that piece if it exists. If all \(E_j\leq T\), the original manifold is
standard. In particular, if \(M\) is nonstandard, the unique piece
\(s\) with \(E_s>T\) satisfies
\begin{equation}\label{eq:cut-discarded}
 \sum_{j\ne s}m_j\leq E-I^+(M)+\gamma_s .
\end{equation}
\end{proposition}

\begin{proof}
Each middle sphere separates, since otherwise a loop could cross it
once, contrary to \(H_1(M;\mathbb Z)=0\). A cut in a current piece
splits it in two: separation holds in the full manifold, and the
connected two-sided collar provides exactly two adjacent components.
The incidence graph of the pieces and cuts is therefore a tree.
Van Kampen and Mayer--Vietoris along its edges, and along the cap
boundaries, give
\[
 \pi_1(M)\cong *_{j}\pi_1(\widehat M_j),\qquad
 H_2(M;\mathbb Z)\cong\bigoplus_j H_2(\widehat M_j;\mathbb Z).
\]
All the factors vanish. By Poincar\'e duality the capped pieces are
simply connected integral homology four-spheres. Hurewicz gives a
map from \(S^4\) representing a fundamental class of each, and the
homological Whitehead theorem makes each a homotopy sphere.

Each cap is an actual positive class on its capped homotopy sphere.
Proposition~\ref{prop:recognition} therefore identifies every piece of
energy \(E_j\leq T\) with \(S^4\). Two energies strictly above \(T\)
would have sum greater than \(2T\), contradicting the total bound.
If two energies equal \(T\), both pieces are standard and every
remaining energy is zero.

Removing the actual cap balls and reusing the original boundary
identifications recovers \(M\). In a standard capped leaf, Palais's
theorem moves the embedded cap ball to a coordinate hemisphere in
the standard sphere. The closed complement of the cap interior is
therefore a standard ball
\cite[Theorem B]{Palais1960Extending}. The resulting boundary
identification extends over a standard ball, since a diffeomorphism
of \(S^3\) is isotopic to an orthogonal map of the same orientation
sign. Use the isotopy in a collar and the orthogonal map towards
the center \cite[p.~553]{Hatcher1983Smale}. Filling the neighboring
piece by the raw leaf or by the cap thus gives the same oriented
diffeomorphism type, by a diffeomorphism that is the identity on
the neighboring piece outside the gluing collar. Root the tree at
the exceptional piece, or at any piece if all are standard.
Successively removing standard leaves identifies \(M\) with the
capped root. In the nonstandard case its metric is an admissible
competitor for \(I^+(M)\), so
\(I^+(M)\leq E_s\leq m_s+\gamma_s\). This gives
\eqref{eq:cut-discarded}.
\end{proof}

The same finite cut proof works with \(T\) replaced throughout
Proposition~\ref{prop:cut-energy} by \(T+\varepsilon_*\), using
Theorem~\ref{thm:recognition-gap} on each positive cap.
In particular the explicit inclusive total \(2T=8\pi^2\) always
suffices once every cap is positive. The scalar-flat wall for a
joined path on the original type still requires a strict subcritical
bound, as in \eqref{eq:zero-wall}.


For one cut there is a useful direct case. If both capped classes
are positive and both raw energies differ from \(E/2\) by at most
\(a\), then \(E/2+a+\max_j\gamma_j\le T\) recognizes both caps and
hence \(M\). For more uneven divisions, the minimizing inequality
\eqref{eq:cut-discarded} brings even a small discarded mass into play.

\begin{corollary}[One-cut nonsplitting]\label{thm:cut}
On a fixed nonstandard type let positive classes have energies
\(e_i\to I^+(M)<2T\). Suppose a marked sphere has full raw sides
\(P_{i,\pm}\) with positive smooth caps of energies at most
\(E_+(g_i;P_{i,\pm})+\gamma_{i,\pm}\). If
\(e_i+\gamma_{i,+}+\gamma_{i,-}\le2T\) at index \(i\), then
\begin{equation}\label{eq:cut-min-bound}
 \min_{\sigma=\pm} E_+(g_i;P_{i,\sigma})
 \le e_i-I^+(M)+\max_{\sigma=\pm}\gamma_{i,\sigma}.
\end{equation}
\end{corollary}
\begin{proof}
In \eqref{eq:cut-discarded} the energy of the discarded side is
bounded by the minimizing excess plus the surcharge on the surviving
side. This proves the inequality, even when the survivor varies with \(i\).
\end{proof}

A short first-order comparison in Proposition~\ref{prop:short-cut}
constructs positive caps by keeping the joined metric strictly below
\(2T\) along a path on \(M\). The inclusive bound above applies once
the capped classes are known to be positive. For a long buffered product
collar the neck paper gives the individual cap estimates directly,
including the localization loss for unequal buffer lengths
\cite{Baker2026MarkedNecks}.

\subsection{Positive entry from the retained exterior}

There is a separate test when a metric has been constructed on a punctured
smooth type and the Yamabe sign of a closed extension is unknown.
Let \(P\) be the compact complement in \(M\) of the interiors of finitely
many disjoint parametrized standard four-balls. Suppose \(h\) is a
smooth metric on \(P\) equal to \(dt^2+\gamma\) on the marked boundary
collars, with \(\gamma\) the unit round three-sphere metric. Write
\(Q_h(v)=\int_P(6|dv|^2+R_hv^2)\,dV_h\).

\begin{proposition}[Round boundary form]\label{prop:round-boundary-form}
There is a smooth closed extension class on \(M\) of energy
\(E_+(h;P)\). It has positive Yamabe constant exactly when every
nonzero \(v\in H^1(P)\), with boundary traces \(f_j\), satisfies
\[
 Q_h(v)+6\sum_j\langle f_j,\sqrt{1-\Delta_\gamma}\,f_j\rangle>0 ,
\]
where \(-\Delta_\gamma\) is nonnegative and the pairing is between
\(H^{1/2}\) and \(H^{-1/2}\), using \(dV_\gamma\). A sufficient condition is positivity of the lowest
eigenvalue of \(-6\Delta_h+R_h\) with
\(\partial_\nu v+v=0\) on each boundary, where \(\nu\) points out of \(P\).
For nonempty boundary this sufficient condition holds whenever
\(R_h\ge0\): vanishing of the Robin form would force a constant function
with zero trace on the connected exterior. Under either sign test,
\(E_+(h;P)\le4\pi^2\) recognizes \(M\), while
\(E_+(h;P)<8\pi^2\) supplies positive subcritical entry.
\end{proposition}

\begin{proof}
Attach unit half-cylinders pointing out from \(P\). Compactify them
in the given markings using a conformal factor equal to one near
\(P\) and to \(e^{-t}\) for large \(t\). In the coordinate
\(r=e^{-t}\) the compactified metric is exactly flat near each added
point. The caps have zero Weyl curvature. On the cylinder the scalar
curvature is six. A spherical harmonic of degree \(l\) has eigenvalue
\(l(l+2)\) for \(-\Delta_\gamma\); this follows by restricting a
homogeneous harmonic polynomial in four variables to the unit sphere.
The least cylinder energy extending its boundary coefficient \(f_l\) is
\[
 6\inf_{w(0)=f_l}\int_0^\infty
       (|w'|^2+(l+1)^2|w|^2)\,dt
   =6(l+1)|f_l|^2 ,
\]
by completing the square; the minimizer is \(f_l e^{-(l+1)t}\).
Orthogonal summation gives the stated boundary form on \(H^{1/2}\)
traces. To justify the energy spaces, let \(X\) denote the exterior
with cylinders, \(f\) the compactifying factor and \(\widehat g=f^2g_X\).
On compactly supported tests the map \(v\mapsto u=fv\) preserves the
Yamabe forms by conformal covariance \cite[\S2]{LeeParker1987Yamabe}.
Adding \(C\|v\|_{2,\widehat g}^2\) adds
\(C\int_X f^2u^2\,dV_{g_X}\) to the cylinder-side form.
For \(C\) sufficiently large the resulting norms are equivalent to the
respective \(H^1\) norms: on \(P\) the factor is one and the scalar is
bounded, and on the cylinders the scalar is six.
Compactly supported tests are dense on \(X\), and tests avoiding the
finitely many added points are dense on the smooth closed metric by
zero point capacity. Thus the correspondence holds on the completions.

Consequently the full quadratic form is the displayed expression
plus a nonnegative sum of zero-trace energies, strictly positive for
each nonzero zero-trace remainder. A nonpositive boundary-form test
extends by its minimizing cap functions to a nonpositive closed test.
On the smooth closed metric the first eigenvalue is attained, so
strict positivity on all nonzero tests is equivalent to positive
Yamabe sign. Finally the eigenvalues \(l+1\) of the boundary operator
are at least one, giving the Robin test by its Rayleigh quotient.
\end{proof}

For a direct recognition test let \(L_h\) be the largest eigenvalue
of either Weyl block on this exterior. If for every nonzero \(v\in H^1(P)\)
\begin{equation}\label{eq:modified-boundary}
 Q_h(v)-6\int_P L_h v^2\,dV_h
      +6\sum_j\langle f_j,\sqrt{1-\Delta_\gamma}\,f_j\rangle>0,
\end{equation}
then \(M\) is standard. Indeed \(L\) has conformal weight minus two
and vanishes on each cap. The preceding extension argument therefore
tests the closed form with potential \(R-6L\) exactly. Its positive
first eigenvalue and bounded potential give coercivity in \(H^1\),
so its modified Yamabe quotient has a positive infimum by Sobolev
embedding. Chen and Zhu's classification and simple connectivity give
the conclusion \cite[Theorem 1.7]{ChenZhu2014Conformal}.
This test uses only the marked exterior. The descent paper gives explicit
round-ended cylinders satisfying it with unbounded exterior Weyl energy
\cite{Baker2026FixedDescent}.

\paragraph{Transfer to an isometric regluing.}
Let \(W\) be a compact oriented smooth contractible four-manifold with
connected boundary \((Y,\gamma)\). Give \(W\) one smooth metric with a
product germ on this actual face. The reflected identity double and
every double glued by an orientation-preserving \(\gamma\)-isometry
are smooth homotopy spheres. If the identity double has positive
Yamabe sign, the free-boundary scalar form on \(W\) is strictly
positive on every nonzero test, by even reflection. Restriction and
addition then give positive sign on every indicated double. Each has
the same chiral energy, namely the sum of the two half energies in
the orientation of \(W\). They also share \(\int L^2dV\), since
orientation reversal interchanges the Weyl blocks. Thus one reflected
metric satisfying the spectral recognition test of
Theorem~\ref{thm:recognition-gap} recognizes all these marked
regluings, including a boundary map that need not extend over \(W\).
A reflected energy below \(2T\) instead gives entry on every such
double. The construction paper derives the free-trace test for a
finite strip and the signature cost at its actual face
\cite[section on short seams]{Baker2026ImmersedHandles}.


The product requirement can be obtained by approximation on the
given half. Suppose a smooth metric \(h\) on \(W\) has positive scalar
curvature, totally geodesic boundary, and induced boundary metric in
the prescribed class \([\gamma]\) in the actual marking. Write
\(B_h=E_+(h;W)+E_-(h;W)\) and \(G_h=\int_W L_h^2dV_h\).
Conformally match \(\gamma\) with zero normal derivative of the
multiplier. In normal coordinates the tangential metric then differs
from its boundary value by \(O(t^2)\). Cutting off this difference
on a layer of width \(b\to0\) gives an exact product with uniformly
bounded curvature and vanishing change in both curvature integrals.
The free scalar form was coercive by positive scalar curvature and
zero mean curvature. The conformal boundary change preserves its
sign, and the layer changes it by \(o(1)\) in the \(H^1\) form norm.
Each compatible isometric double consequently has positive classes
with \(e\to B_h\) and \(P\to4B_h+6G_h\)
\cite[section on the eta-zero filling]{Baker2026ImmersedHandles}.
Either \(B_h\le T+\varepsilon_*\) or \(4B_h+6G_h\le8T+\eta\)
recognizes every one of these fixed smooth types, including equality
with the approximating classes approaching from above.

For the indicated Mazur half there are smooth positive-scalar-curvature
extensions \(g\) of the exact eta-zero boundary metric. The extra
condition needed on some such extension can be tested entirely on
boundary traces. Let \(\mathcal N_g\) be the outward normal derivative
of the \((-6\Delta_g+R_g)\)-harmonic Dirichlet extension and \(H_g\)
the outward mean curvature in the trace convention. A positive
totally geodesic filling in the prescribed class exists exactly when,
for some extension, \(6\mathcal N_g+2H_g\) has positive quadratic
form on all nonzero traces. The positive first Robin eigenfunction
gives a conformal metric with minimal boundary; boundary deformation
then makes it totally geodesic with the same induced metric.
In the converse direction one fixes the target boundary metric by a
conformal multiplier and solves a positive Dirichlet problem to restore
positive scalar curvature. The sign for one actual extension and the
half-energy budget on the resulting metric remain to be established
\cite[section on filling the marked boundary class]{Baker2026ImmersedHandles}.
